%%%%%%%%%%%%%%%%%%%%%%%%%%%%%%%%%%%%%%%%%%%%%%%%%%%%%%%%%%%
% --------------------------------------------------------
% Tau
% LaTeX Template
% Version 2.4.3 (01/09/2024)
%
% Author: 
% Guillermo Jimenez (memo.notess1@gmail.com)
% 
% License:
% Creative Commons CC BY 4.0
% --------------------------------------------------------
%%%%%%%%%%%%%%%%%%%%%%%%%%%%%%%%%%%%%%%%%%%%%%%%%%%%%%%%%%%

\documentclass[9pt,a4paper,twoside]{tau-class/tau}
\usepackage{caption}
\usepackage[utf8]{inputenc}
\usepackage[dvipsnames]{xcolor}
\usepackage{xcolor}
\usepackage{tikz}
\usetikzlibrary{positioning,arrows.meta,shadows,shapes.geometric}
\captionsetup[table]{position=bottom}
\usepackage{tikz}
\usepackage{pgfplots}
\pgfplotsset{compat=1.17}
\usepackage{booktabs}
\usepackage{amsmath}
\usepackage{siunitx}
\usepackage{tcolorbox}
\usepackage{graphicx}
\usetikzlibrary{fit,shadows}
\usetikzlibrary{shapes.geometric,shapes.misc}

%----------------------------------------------------------
% TITLE
%----------------------------------------------------------

\journalname{CoMe Project Proposal --- Computational Mechanics}
\title{Buckling Analysis of Cylindrical Shells using\\
Bogner-Fox-Schmit (BFS) Elements}

%----------------------------------------------------------
% AUTHORS, AFFILIATIONS AND PROFESSOR
%----------------------------------------------------------

\author[a,1]{Siddharthan Somasundaram (Matriculation No.: 254304)}

%----------------------------------------------------------

\affil[a]{Otto-von-Guericke-Universit\"{a}t Magdeburg, Germany}

\professor{Supervisor: Dr.-Ing.\ Sascha Eisentr\"{a}ger \quad
Institute of Materials, Technologies and Mechanics --- Computational Mechanics \quad
Prof.\ Dr.-Ing.\ Daniel Juhre}

%----------------------------------------------------------
% FOOTER INFORMATION
%----------------------------------------------------------

\institution{Otto-von-Guericke-Universit\"{a}t Magdeburg, Germany}
\footinfo{Version 1.0}
\theday{\today}
\leadauthor{Siddharthan Somasundaram}
\course{CoMe Project --- Computational Mechanics (2025)}

%----------------------------------------------------------
% ABSTRACT AND KEYWORDS
%----------------------------------------------------------

\begin{abstract}
This project proposes the implementation and validation of a linear buckling
analysis framework for cylindrical shells using the classical
Bogner-Fox-Schmit (BFS) finite element. BFS elements are $C^1$-continuous
plate and shell elements constructed from bi-cubic Hermitian polynomials,
making them well-suited for thin-shell stability problems. Building upon an
existing validated Matlab code base for BFS-based flat plate analysis, this
work extends the formulation to cylindrical shell geometry through three
structured phases: the derivation of the cylindrical shell element with
48~degrees of freedom per element; the implementation of a pre-buckling
in-plane stress solver; and the assembly of the geometric stiffness matrix
required for eigenvalue-based buckling analysis. The project targets a
simply-supported cylindrical shell under uniform axial compression and
validates the computed critical buckling load against the classical Donnell
analytical solution and Ansys reference results. This work constitutes the
direct computational foundation for the subsequent Master's thesis, which
will extend the methodology to cylindrical shells with multiple cutouts
using the Cut-BFS (CutBFS) approach based on the Finite Cell Method.
\end{abstract}

\keywords{Bogner-Fox-Schmit elements, cylindrical shell, buckling analysis,
geometric stiffness, Finite Cell Method, CutBFS}

%----------------------------------------------------------

\begin{document}
		
    \maketitle 
    \thispagestyle{firststyle}
    \tauabstract 
    \tableofcontents
    \linenumbers 

%----------------------------------------------------------
\section{Introductory Paragraph}
%----------------------------------------------------------

\subsection{Problem Statement and Importance}

\taustart{T}hin-walled cylindrical shell structures are fundamental
load-bearing components in modern engineering applications, including
aerospace fuselages, rocket bodies, submarine pressure hulls, and offshore
pipelines. Under axial compressive loading, these structures are susceptible
to a sudden and catastrophic loss of stability known as \emph{buckling} ---
a failure mode that is particularly dangerous because it occurs abruptly,
often at load levels well below the material yield strength. The accurate
prediction of the critical buckling load and the associated deformation mode
shape is therefore an essential requirement in the structural design and
certification of such components~\cite{timoshenko1961,abaqus2009}.

The Bogner-Fox-Schmit (BFS) element, introduced in 1965, is a
high-efficiency $C^1$-continuous finite element based on a bi-cubic
Hermitian polynomial approximation~\cite{bogner1965}. This continuity
property makes it particularly well-suited for the analysis of thin shells
under bending and buckling, as it analytically satisfies slope continuity
between adjacent elements without penalty or Lagrange multiplier techniques.
Despite its proven accuracy and implementation simplicity, the BFS element
has seen limited application in modern structural analysis, primarily because
its requirement for a structured rectangular mesh restricts its direct use on
geometrically complex structures. Extending the BFS element to cylindrical
shell geometries represents a natural and important step towards broader
application in engineering practice~\cite{eisentraeger2022}.

\subsection{Scope of the Project}

This project implements and validates a linear buckling analysis framework
for cylindrical shells using BFS elements. The work builds directly upon an
existing validated Matlab code base for BFS-based plate analysis
\cite{eisentraeger2022} and extends it to curved shell geometry through
three structured phases:
%
\begin{enumerate}
    \item Formulation and implementation of the cylindrical BFS shell element
          with 48 degrees of freedom per element (12 per node, accounting for
          three displacement components);
    \item Development of the pre-buckling in-plane (membrane) stress solver;
    \item Construction and assembly of the geometric stiffness matrix
          $\mathbf{K}_G$ required for eigenvalue-based buckling analysis.
\end{enumerate}
%
The project considers a simply-supported cylindrical shell under uniform
axial compression --- the canonical benchmark for shell buckling --- and
validates results against the classical Donnell analytical solution. This
work constitutes the direct computational foundation for the subsequent
Master's thesis, which will extend the methodology to cylindrical shells
with multiple cutouts using the CutBFS approach.

%----------------------------------------------------------
\section{Brief State of the Art}
%----------------------------------------------------------

\subsection{Review of Relevant Literature}

Linear buckling analysis of cylindrical shells has a long history in
structural mechanics. The classical analytical solution for a
simply-supported cylindrical shell under axial compression, derived from
Donnell's shallow shell equations, provides the reference critical stress
\cite{timoshenko1961}:
%
\begin{equation}
    \sigma_{\mathrm{cr}} 
    = \frac{E}{\sqrt{3\left(1-\nu^2\right)}} \cdot \frac{h}{R},
    \label{eq:donnell}
\end{equation}
%
where $E$ is Young's modulus, $\nu$ is Poisson's ratio, $h$ is the wall
thickness, and $R$ is the cylinder radius. This expression remains the
primary benchmark for numerical methods. Finite element analysis of shell
buckling has advanced significantly, with commercial codes such as Abaqus and
Ansys offering robust shell element formulations validated against
experimental data~\cite{abaqus2009}.

The original BFS element formulation for cylindrical shells was proposed by
Bogner, Fox and Schmit themselves~\cite{bogner1965}, demonstrating the
element's suitability for curved geometries. More recently,
Eisentr\"{a}ger~et~al.~\cite{eisentraeger2022} demonstrated that BFS
elements, when combined with the Finite Cell Method (FCM) in the so-called
CutBFS approach, can accurately predict the critical buckling load of
rectangular plates with multiple cutouts. The CutBFS method replaces the
geometry-conforming mesh requirement with a structured Cartesian background
mesh, thereby retaining the accuracy and simplicity of BFS elements while
enabling the analysis of geometrically complex structures. In a complementary
development, Burman~et~al.~\cite{burman2020} established the mathematical
foundation for cut BFS elements and proved optimal convergence rates in the
context of plate bending. Furthermore, the Shell Finite Cell Method
\cite{shellFCM2011} demonstrated that fictitious domain approaches can be
generalised to thin-walled shell structures using high-order $p$-version
elements; however, no equivalent framework based on the $C^1$-continuous BFS
element has been developed for cylindrical shell geometries.

\subsection{Identified Gap}

The central gap in the existing literature is the absence of a validated,
open-source BFS-based buckling analysis framework for cylindrical shells.
While BFS plate buckling has been comprehensively studied~\cite{eisentraeger2022},
and while cylindrical shell buckling has been analysed using alternative finite
element formulations~\cite{timoshenko1961,abaqus2009}, the combination of
BFS elements with cylindrical shell geometry and buckling analysis has not
been implemented or validated in a modern Matlab environment. This project
directly addresses this gap and provides the validated shell code base
required for the CutBFS cylinder extension in the Master's thesis.

%----------------------------------------------------------
\section{Open Research Question}
%----------------------------------------------------------

The central research question of this project is:

\begin{tcolorbox}[colback=gray!10, colframe=gray!50, rounded corners]
\textit{Can BFS elements, implemented within the existing CutBFS Matlab
framework, accurately predict the critical buckling load and mode shape of a
simply-supported cylindrical shell under axial compression, and be validated
against the classical Donnell analytical solution?}
\end{tcolorbox}

\noindent This overarching question is addressed through three specific and
graduated sub-questions:

\begin{enumerate}
    \item \textbf{Pre-buckling stress:} What is the in-plane membrane stress
    distribution in the cylindrical shell wall under a unit axial compressive
    load, and does the BFS-based plane stress solver reproduce the expected
    uniform stress state?

    \item \textbf{Critical load:} What is the critical load factor $\lambda$
    predicted by the BFS eigenvalue analysis, and how does it compare to the
    Donnell analytical solution given in Equation~\eqref{eq:donnell}?

    \item \textbf{Mode shape:} What is the first buckling mode shape (number
    of axial half-waves $m$ and circumferential waves $n$), and does it match
    the analytically predicted wave pattern?
\end{enumerate}

%----------------------------------------------------------
\section{Methodology}
%----------------------------------------------------------

\subsection{Objectives and Techniques}

The project is structured into four implementation phases, each building
directly on the previous:

\paragraph{Phase 1 --- Cylinder geometry and mesh.}
The cylindrical shell of radius $R$, length $L$, and wall thickness $h$ is
\emph{unfolded} into a rectangular domain using the coordinate transformation
$x$ (axial direction) and $y = R\theta$ (circumferential direction). A
structured rectangular mesh is generated, enabling direct reuse of the
existing BFS element infrastructure. The 48-DOF per element formulation is
adopted, with each of the three displacement components ($u$, $v$, $w$)
carrying four Hermitian DOFs per node:
%
\begin{equation}
    f_i,\quad
    \left(\frac{\partial f}{\partial x}\right)_i,\quad
    \left(\frac{\partial f}{\partial y}\right)_i,\quad
    \left(\frac{\partial^2 f}{\partial x\,\partial y}\right)_i,
    \quad f \in \{u,\,v,\,w\}.
    \label{eq:dofs}
\end{equation}

\paragraph{Phase 2 --- Pre-buckling in-plane stress analysis.}
A 32-DOF plane stress stiffness matrix $\mathbf{K}_{\mathrm{plane}}$ is
assembled using the BFS shape functions for the in-plane displacement
components $u$ and $v$. A unit axial compressive load is applied as
energetically equivalent nodal forces at the cylinder ends. The resulting
membrane stress resultants $(\tilde{N}_x,\,\tilde{N}_y,\,\tilde{N}_{xy})$
are computed at each integration point and stored for Phase~3:
%
\begin{equation}
    \mathbf{K}_{\mathrm{plane}}\,\mathbf{U}_{\mathrm{ref}}
    = \mathbf{F}_{\mathrm{edge}}.
    \label{eq:prestress}
\end{equation}

\paragraph{Phase 3 --- Geometric stiffness and buckling eigenvalue.}
The geometric (initial-stress) stiffness matrix $\mathbf{K}_G$ is assembled
by integrating the product of the BFS shape function gradients and the
pre-buckling membrane stress tensor over each element:
%
\begin{equation}
    \mathbf{k}_G^{(e)} 
    = \int_{-1}^{1}\int_{-1}^{1}
      \mathbf{L}_g^\top\,\boldsymbol{\Sigma}\,\mathbf{L}_g\,
      \det(\mathbf{J})\,\mathrm{d}\xi\,\mathrm{d}\eta,
    \label{eq:KG}
\end{equation}
%
where $\mathbf{L}_g$ is the gradient operator applied to the BFS shape
functions and $\boldsymbol{\Sigma}$ is the pre-buckling stress matrix. The
generalised eigenvalue problem
%
\begin{equation}
    \left(\mathbf{K}_{\mathrm{bending}} + \lambda\,\mathbf{K}_G\right)
    \boldsymbol{\Phi} = \mathbf{0}
    \label{eq:eig}
\end{equation}
%
is solved using Matlab's \texttt{eigs} function to extract the critical load
factor $\lambda$ and the corresponding mode shape $\boldsymbol{\Phi}$.

\paragraph{Phase 4 --- Validation and post-processing.}
The computed critical load is compared against the Donnell analytical
solution (Eq.~\eqref{eq:donnell}). Convergence with respect to mesh
refinement is assessed. The buckling mode shape is visualised using
Matlab's \texttt{surf} command and compared qualitatively against the
analytically predicted diamond patterns.

\subsection{Simulation Tools}

All implementation is carried out in Matlab, extending the existing BFS
plate code provided by the research group~\cite{eisentraeger2022}. No
commercial finite element software is used for the primary analysis; Ansys
Mechanical is employed solely as an independent reference for
cross-validation of the critical buckling load.

\subsection{Anticipated Limitations}

The analysis is restricted to \emph{linear} (classical) buckling, which
neglects geometric nonlinearities and initial imperfections known to
significantly reduce the actual buckling load of cylindrical shells in
practice. The simply-supported boundary condition is assumed throughout;
clamped or free-edge conditions are not considered in this project but are
identified as natural extensions. Cutouts are explicitly excluded from this
project scope and are reserved for the Master's thesis.

%----------------------------------------------------------
\section{Outcome}
%----------------------------------------------------------

\subsection{Planned Results}

The expected technical outcomes of this project are:
%
\begin{itemize}
    \item A validated Matlab implementation of the 48-DOF cylindrical BFS
          shell element, including the pre-buckling plane stress solver and
          the geometric stiffness matrix $\mathbf{K}_G$.
    \item A convergence study demonstrating that the BFS-based critical
          buckling load converges to the Donnell analytical solution with
          mesh refinement, with a target accuracy of less than~5\%.
    \item Contour plots of the pre-buckling membrane stress distribution,
          confirming the expected uniform axial stress state under end
          compression.
    \item Visualisation of the first buckling mode shape on the cylindrical
          surface, showing the characteristic diamond pattern with the
          analytically predicted wave numbers $m$ and $n$.
    \item A parametric study of the effect of the radius-to-thickness ratio
          $R/h$ on the critical buckling load, comparing BFS results
          against the analytical formula and Ansys reference solutions.
\end{itemize}

\subsection{Contribution to Knowledge}

This project provides the first validated implementation of BFS shell
elements for cylindrical shell buckling analysis within the existing CutBFS
Matlab framework. It establishes a tested and documented code base that
directly enables the Master's thesis extension to cylindrical shells with
multiple cutouts using the CutBFS method. The parametric validation against
analytical and commercial FEM solutions provides confidence in the
correctness of the implementation and quantifies the accuracy of the BFS
element for this class of problem.

%----------------------------------------------------------
\section{Conclusion}
%----------------------------------------------------------

This project proposes the implementation and validation of a linear buckling
analysis framework for cylindrical shells using Bogner-Fox-Schmit elements,
extending an existing validated Matlab code base for BFS plate analysis. The
work is strategically scoped as the direct computational foundation for the
subsequent Master's thesis on CutBFS-based buckling analysis of cylindrical
shells with multiple cutouts.

By completing this project, the key missing components --- the 48-DOF
cylindrical shell element, the pre-buckling plane stress solver, and the
geometric stiffness matrix $\mathbf{K}_G$ --- will be implemented and
validated, leaving the thesis to focus on the novel CutBFS extension with
cutout handling and parametric optimisation. The immediate next step is a
detailed study of the existing plate code structure alongside the
mathematical roadmap document provided by the supervisor, before
implementation begins in Matlab.

%----------------------------------------------------------
% BIBLIOGRAPHY
%----------------------------------------------------------

\begin{thebibliography}{9}

\bibitem{timoshenko1961}
S.~Timoshenko, J.~Gere,
\textit{Theory of Elastic Stability}, 2nd~ed.
McGraw-Hill, New York, 1961.

\bibitem{abaqus2009}
Dassault Syst\`{e}mes Simulia Corp.,
\textit{Abaqus Theory Manual}, vol.~6.6, sec.~3.6.4.
Providence, RI, USA, 2009.

\bibitem{bogner1965}
F.~K.~Bogner, R.~L.~Fox, L.~A.~Schmit,
``The generation of interelement compatible stiffness and mass matrices
by the use of interpolation formulae,''
in \textit{Proc.\ 1st Conf.\ Matrix Methods in Structural Mechanics},
Wright Patterson AF Base, Ohio, 1965, pp.~397--444.

\bibitem{eisentraeger2022}
S.~Eisentr\"{a}ger, J.~Kiendl, G.~Michaloudis, R.~Duy, Y.~Vetyukov,
``Stability analysis of plates using cut Bogner-Fox-Schmit elements,''
\textit{Computers \& Structures}, vol.~270, p.~106854, 2022.
\texttt{doi:10.1016/j.compstruc.2022.106854}

\bibitem{burman2020}
E.~Burman, P.~Hansbo, M.~G.~Larson,
``Cut Bogner-Fox-Schmit elements for plates,''
\textit{Adv.\ Model.\ Simul.\ Eng.\ Sci.}, vol.~7, no.~27, 2020.
\texttt{doi:10.1186/s40323-020-00164-3}

\bibitem{shellFCM2011}
D.~Schillinger, M.~Ruess, N.~Zander, Y.~Bazilevs, A.~D\"{u}ster, E.~Rank,
``Small and large deformation analysis with the $p$- and B-spline versions
of the Finite Cell Method,''
\textit{Comput.\ Methods Appl.\ Mech.\ Eng.}, vol.~211, 2011.

\bibitem{duy2021}
R.~Duy,
\textit{Fictitious Domain Approach for Optimizing Stability Boundaries of
Plates with Cutouts},
Master's thesis, TU~Wien, Vienna, Austria, 2021.

\end{thebibliography}

\end{document}